\section{Next-Edge Reduction and Cache-Amortized Regret}
\label{app:next}
\subsection{Signed Remainder Comparison}
Omit time indices and the common cost $h$. Write $a$ for the selected action and $a^*$ for the full-cost comparator. On a reward-confidence round,
\begin{align*}
\mu_{a^*}-\lambda e_{a^*}
&\le\widehat Q(a^*)+b(a^*)-\lambda e_{a^*}\\
&\le\widehat Q(a)+b(a)-\lambda e_a\\
&\le\mu_a-\lambda e_a+2b(a).
\end{align*}
Subtract $\lambda r_{a^*}$ on the left and express the right relative to $\mu_a-\lambda(e_a+r_a)$. This gives~\eqref{eq:signedtail}. If $r_{a^*}\ge0$ and $\lambda\ge0$, drop $-\lambda r_{a^*}$. A common remainder cancels exactly before applying confidence bounds, so the complete-cost theorem with constant $8$ applies in that reduction. Common remainders can vary with time and cache state without changing this conclusion.

For nonnegative unequal remainders, use $E_t=\lambda r_t(a_t)$ in the generic counting argument. If $E_t>\Delta_t/2$, charge the gap to $2E_t$. Otherwise a positive gap requires $N_t(a_t)\le32\log T/\delta^2$. Initialization contributes at most $KD$, counts below the threshold contribute at most $KD(32\log T/\delta^2+1)$, and the expected failed-confidence contribution is at most $4KD$. Integration proves~\eqref{eq:nextbound}. The Lean theorem \leanname{expected_next_edge_regret} composes this argument with the reward-stream MGF bound; it does not assume a successful confidence event.

\subsection{From Physical Work to a Forecast Budget}
The forecast in~\eqref{eq:tailforecast} concerns the selected action under the actual completion and panel law. Its conditioning information includes that choice; the bill $X_t$ is the subsequent realized omitted cold work. For each path, persistent materialization gives
\begin{equation}
\begin{split}
\sum_{t<T}X_t
&\le\sum_{t<T}\sum_{r\in R_t\setminus M_t}k_r\\
&=W(M_T)-W(M_0)\le B_{\mathrm{cold}}.
\end{split}
\end{equation}
No independence between different resource demands is used. For each $t$, the conditional-expectation identity gives $\E r_t(a_t)=\E X_t$. Finite sums commute with integration, hence $\E\sum_t r_t(a_t)\le B_{\mathrm{cold}}$. The formal proof uses mathlib's conditional expectation and \leanname{integral_condExp}; \leanname{expected_forecast_cold_budget} combines it with the actual ledger bound. Finally, \leanname{expected_next_edge_regret_with_cache} supplies the complete selection-to-regret composition in Theorem~\ref{thm:nextcache}.

The finite constant is attached to a fixed physical resource universe. A growing query pool, new execution namespaces, eviction, or repeated stochastic re-execution can create new chargeable work. Such work must be represented by new resources or separate overhead, and its cumulative contribution can grow with $T$. A fixed biased forecast is not made conditionally correct by a finite physical cache.

The theorem is local. Summing node-level regrets can charge a rollout's omitted work at several nodes, so a whole-tree bound must retain that multiplicity and any reward-drift terms. Physical billing still charges the execution once.

\subsection{The Implemented Immediate Penalty}
Let $B_t=\widehat C_t(s)$ denote the current prefix's predicted missing bill. The implementation computes
\begin{equation}
\widehat e_t(s,a)=\max\{0,\widehat C_t(sa)-B_t\}.
\label{eq:immediateimpl}
\end{equation}
If $\widehat C_t(sa)=0$ and $B_t\ge0$, this value is zero. If it is zero for every candidate, subtracting its penalty leaves the entire UCB maximizer set unchanged. These are \leanname{cached_child_next_penalty_zero} and \leanname{zero_next_penalty_recovers_ucb}.

Equation~\eqref{eq:immediateimpl} is an ideal next-edge acquisition cost only when its preview corresponds to work that the eventual evaluation actually executes. A sparse deeper checkpoint can bypass that edge and make a complete bill smaller than a partial-prefix preview. Then $c_t-e_t$ can be negative, violating the nonnegative-remainder assumption. The signed comparison remains algebraically valid, but the one-sided cache-budget theorem needs a valid physical decomposition. Forecast calibration and a mismatched reference law are additional errors, as quantified in Proposition~\ref{prop:costerror} and Theorem~\ref{thm:errors}.

\subsection{Exact Edges with Recurring Omitted Work}
In Theorem~\ref{thm:nextlinear}, $e_A+r_A=e_B+r_B=0.4$, with both remainders nonnegative. If $n_A+n_B=T$, then $R_T^{\mathrm{full}}=0.1n_B$ and $R_T^{\mathrm{next}}=0.2n_A$, giving the identity used in its proof. The two initialization pulls and all possible subsequent tie-breaking choices are covered by \leanname{next_edge_linear_regret}. This construction can represent repeated evaluation or repair work that is not memoizable. It does not contradict Theorem~\ref{thm:nextcache}, because the omitted recurring bills do not obey a fixed finite materialization budget.

\subsection{Persistent Forecast Bias}
\begin{theorem}[Linear true regret from biased cost]
\label{thm:bias}
Consider deterministic rewards $\mu_A=0.8$, $\mu_B=0.7$, true costs $c_A=c_B=0.1$, predicted costs $\widehat c_A=0.4$, $\widehat c_B=0.1$, and $\lambda=1$. Initialize $A$ and then $B$, and thereafter maximize the plug-in UCB index with exact reward means. For every such trajectory and $T\ge1$,
\begin{equation}
R_T^{\mathrm{true}}\ge T/10-40\log T-2/5.
\label{eq:lower}
\end{equation}
\end{theorem}
\begin{proof}
True net means are $0.7,0.6$, whereas proxy net means are $0.4,0.6$. Each choice of $A$ incurs proxy regret $0.2$ and zero true regret; each choice of $B$ incurs true regret $0.1$ and zero proxy regret. Hence $R_T^{\mathrm{true}}=T/10-R_T^{\mathrm{proxy}}/2$. Exact rewards have no failed-confidence rounds. The deterministic bound used in Theorem~\ref{thm:log}, with $K=2$ and $D=\delta=0.2$, gives $R_T^{\mathrm{proxy}}\le80\log T+4/5$. Substitute this inequality.
\end{proof}

The example also fits a setting where execution prefixes are already cached and the remaining bill is unshared evaluation overhead. Caching does not repair a biased estimate of that bill.

